\documentclass{beamer}
\usetheme{Madrid}
\usecolortheme{whale}

\title{Definisi Terkait Teorema 7.5}
\subtitle{Ekspansi Polinomial Karakteristik Matriks Graf}
\author{Analisis Literatur Norman Biggs}
\date{\today}

\begin{document}

\frame{\titlepage}

\begin{frame}
  \frametitle{Objek Utama: Graf dan Matriks Q}
  \begin{block}{Graf Sederhana $\Gamma$}
    Graf $\Gamma$ terdiri dari himpunan titik $V\Gamma$ dan himpunan sisi $E\Gamma$. Hubungan antar keduanya ditentukan oleh relasi insidensi.
  \end{block}
  \begin{block}{Matriks Signless Laplacian}
    Matriks $Q$ didefinisikan sebagai:
    $$Q = A + \Delta$$
    di mana $A$ adalah matriks ketetanggaan dan $\Delta$ adalah matriks diagonal yang entri-entrinya berisi derajat (valency) dari setiap titik di $\Gamma$.
  \end{block}
\end{frame}

\begin{frame}
  \frametitle{Matriks Insidensi $D$}
  Matriks insidensi $D$ berukuran $|V\Gamma| \times |E\Gamma|$ menghubungkan titik dengan sisi.
  \begin{exampleblock}{Identitas Fundamental}
    $$Q = DD^t$$
  \end{exampleblock}
  Identitas ini menjadi jembatan utama antara sifat aljabar matriks $Q$ dengan struktur topologis dari graf $\Gamma$.
\end{frame}

\begin{frame}
  \frametitle{Polinomial Karakteristik $\tau(\Gamma; \lambda)$}
  Fungsi ekspansi didefinisikan sebagai:
  $$\tau(\Gamma; \lambda) = \det(\lambda I + Q) = \lambda^n + q_1 \lambda^{n-1} + \dots + q_{n-1}\lambda + q_n$$
  \begin{alertblock}{Definisi Koefisien $q_i$}
    Setiap koefisien $q_i$ ($1 \le i \le n$) merupakan jumlahan dari seluruh minor utama (principal minors) berukuran $i \times i$ dari matriks $Q$.
  \end{alertblock}
\end{frame}

\begin{frame}
  \frametitle{Nilai Spesifik Koefisien $q_i$}
  Beberapa nilai koefisien yang sudah diketahui dari observasi elementer:
  \begin{itemize}
    \item $q_1 = 2|E\Gamma|$ (Dua kali jumlah sisi / jumlah derajat)
    \item $q_{n-1} = n\kappa(\Gamma)$ (Di mana $\kappa(\Gamma)$ adalah jumlah pohon rentang graf)
    \item $q_n = 0$ (Matriks $Q$ selalu bersifat singular)
  \end{itemize}
\end{frame}

\begin{frame}
  \frametitle{Submatriks $D(X, Y)$ \& Teorema Binet-Cauchy}
  Misalkan $X \subseteq V\Gamma$ dan $Y \subseteq E\Gamma$ dengan $|X| = |Y| = i$.
  \begin{itemize}
    \item $D(\cdot, Y)$: Submatriks $D$ yang hanya mengambil kolom dari himpunan sisi $Y$.
    \item $D(X, Y)$: Submatriks dari $D(\cdot, Y)$ yang hanya mengambil baris dari himpunan titik $X$.
  \end{itemize}
  \begin{block}{Ekspansi Binet-Cauchy}
    $$\det Q_X = \sum_{Y} (\det D(X, Y))^2$$
  \end{block}
\end{frame}

\begin{frame}
  \frametitle{Struktur Kombinatorika: Hutan $\Phi$}
  \begin{block}{Hutan (Forest) $\Phi$}
    Hutan $\Phi = \langle Y \rangle$ adalah subgraf dari $\Gamma$ yang tidak memuat siklus dan memiliki tepat $i$ sisi.
  \end{block}
  \begin{block}{Fungsi $p(\Phi)$}
    Banyaknya cara untuk membuang tepat satu titik dari setiap komponen terhubung pada hutan $\Phi$. Jika setiap komponen memiliki jumlah titik $v_1, v_2, \dots, v_k$, maka:
    $$p(\Phi) = v_1 \times v_2 \times \dots \times v_k$$
  \end{block}
\end{frame}

\begin{frame}
  \frametitle{Kesimpulan Teorema 7.5}
  Formulasi akhir untuk koefisien polinomial karakteristik:
  $$q_i = \sum p(\Phi) \quad (1 \le i \le n)$$
  di mana jumlahan dilakukan atas semua subgraf sisi dari $\Gamma$ yang memiliki $i$ sisi dan berbentuk hutan (forest).
\end{frame}

\end{document}